% Source: https://artofproblemsolving.com/community/c3h2795806_2002_midmichigan_mathematical_olympiad__grades_79
% Stable identity: aops-topic-2795806-q01
Number the coins $1$--$12$.

\emph{Weighing 1:} $\{1,2,3,4\}$ against $\{5,6,7,8\}$.

\emph{If it balances,} coins $1$--$8$ are genuine and the fake is among $9$--$12$. Weigh $9$ against genuine coin $1$, then $10$ against $1$, then $11$ against $1$; the first imbalance identifies the fake, and if all three balance the fake is $12$. This uses at most $4$ weighings.

\emph{If it does not balance,} say the left pan is heavier (the other case is symmetric). Then $9$--$12$ are genuine, and either one of $1,2,3,4$ is heavy or one of $5,6,7,8$ is light.

\emph{Weighing 2:} $\{1,2,5\}$ against $\{3,6,9\}$.
\begin{itemize}
\item Left heavier: the fake is $1$ or $2$ (heavy) or $6$ (light). \emph{Weighing 3:} $1$ against $2$; the heavier is fake, and if they balance the fake is $6$.
\item Right heavier: the fake is $3$ (heavy) or $5$ (light). \emph{Weighing 3:} $3$ against $9$; if $3$ is heavier it is fake, otherwise $5$ is.
\item Balance: the fake is $4$ (heavy) or $7$ or $8$ (light). \emph{Weighing 3:} $7$ against $8$; the lighter is fake, and if they balance the fake is $4$.
\end{itemize}
In every case the counterfeit is found within four weighings (in the unbalanced case, within three). $\blacksquare$
